\bookchapter{第三十四章：車機互聯深指南——手機與汽車之融合}

\chapterintro{汽車方自“四個輪子加沙發”進化爲“四個輪子加智能手機”。無論汝開之乃燃油車抑新能源，手機與車機之關係皆決定了汝之駕駛體驗。此一章帶汝釐清車機互聯之三大路線、各派系車機系統之特點，以及彼使開車更順手之小技巧。}

\booksection{一、車機互聯之三大路線}

當前手機與汽車“聯網”主要有三種思路，理解矣彼等遂理解矣整個生態：

\begin{itemize}
  \item \textbf{手機投屏至車機}：車機唯當一塊"顯示器"與"音箱"，運算與地圖數據皆來自手機。代表方案：Apple CarPlay、Google Android Auto、華爲HiCar、百度CarLife。優點乃手機生態直接搬至車上，不用重新裝App；缺點乃車機本身之硬件能力爲閒置。
  \item \textbf{手機與車機深融合}：車機乃獨立系統，但與手機共享賬號、應用、數據，甚至能"接力"任務。代表方案：鴻蒙座艙（問界/智界）、比亞迪DiLink、蔚來NIO OS、理想車機系統。優點乃體驗一體化；缺點乃通常綁定同一品牌生態。
  \item \textbf{CarPlay 2.0深集成版}：2025年蘋果發佈之新一代CarPlay，不再只為"投屏"，而乃接管車機之儀表盤、空調控制、座椅調節等核心功能。法拉利、阿斯頓·馬丁、保時捷等品牌率先搭載，標誌車機系統進入"手機廠商主導"之新階段。
\end{itemize}

\booksection{二、CarPlay / Android Auto / HiCar 對較}

三大主流投屏方案各有側重，選車或選手機時可對照下表：

\begin{center}
\begin{tabularx}{\linewidth}{|X|X|X|X|X|}
\hline
\textbf{方案} & \textbf{支手機系統} & \textbf{支車型} & \textbf{能豐度} & \textbf{國內可用性} \tabularnewline
\hline
Apple CarPlay & 僅 iPhone & 極廣，絕大多數車型支 & 高，地圖/音樂/電話/消息均可 & 好，原生可用 \tabularnewline
Android Auto & Android（國內手機需裝服框架） & 較廣，合資/進口車居多 & 高，與CarPlay相當 & 通常，需科學上網激活 \tabularnewline
華爲HiCar & 華爲手機爲主，部支其他安卓 & 國產車與部合資車型 & 較高，可與鴻蒙座艙聯動 & 甚好，本土生態 \tabularnewline
百度CarLife & iOS 與 Android 均可 & 較廣，舊車型後裝常見 & 中，更新較慢 & 甚好，本土生態 \tabularnewline
\hline
\end{tabularx}
\end{center}

\booksection{三、鴻蒙座艙：手機與車機之無縫流轉}

若曰CarPlay乃“手機投至車機”，鴻蒙座艙則乃“手機與車機本即一家”。華爲手機與問界/智界車機之協同，代表矣當前國內車機互聯之最高水準：

\begin{itemize}
  \item \textbf{超級桌面}：車機直接調用手機上之App，無需於車機重新安裝。手機裏於聽之音樂、於導航之路線，上車自動續接至車機屏幕，下車再回到手機，全程無縫。
  \item \textbf{應用接力}：手機上設好之導航，關車門時車機自動接管；到達停車場，下車後手機繼續步行導航至最終目的地。"終一公里"不再斷檔。
  \item \textbf{手機車機共用賬號}：華爲賬號一登，車機自動同步汝之音樂會員、地圖藏、日程、聯繫人。車機與手機乃同一身份之不同終端，而非兩套獨立系統。
\end{itemize}

\booksection{四、造車新勢力之車機系統}

國產新能源車之車機已自“導航+音樂”進化爲移動智能終端，各家之玩法各有特色：

\begin{itemize}
  \item \textbf{蔚來NIO OS}：NOMI語音助手乃行業標杆，能識別多指令、上下文連續對話，尚帶情感表情。NIO Phone則將手機作成"車鑰匙+第二屏"，深度綁定蔚來生態。
  \item \textbf{理欲車機}：多屏交互乃亮點，副駕屏可獨立觀劇、孩子屏可觀動畫片。"任務大師"使汝自定義場景，比如"開門自動暫停音樂、關門繼續播放"。
  \item \textbf{小米澎湃車機}：與小米手機、手錶、家電深度聯動。手機上設之鬧鐘將於車機提醒；車上導航回家，家裏空調提前開啓。HyperOS打通矣"人車家全生態"。
  \item \textbf{較亞迪DiLink}：標誌性之旋轉屏，豎屏觀短視頻、橫屏觀地圖，自適應App方向。自帶應用市場，生態豐富度於傳統車企中領先。
\end{itemize}

\booksection{五、手機當車機——能用嗎？}

若汝之車機太老或太難用，能否直接用手機“頂上”？答案乃：可，但要觀場景。

\begin{itemize}
  \item \textbf{摩託羅拉Ready For / 桌面模式}：部摩託羅拉、三星、華爲手機支"桌面模式"，連上顯器便乃一臺簡易電腦。理上接車機屏亦能用，但實際兼容性參差不齊，且缺汽車專屬交互（如方向盤按鍵）。
  \item \textbf{手機支架 vs 原車車機}：若原車車機卡頓、地圖不更新，一個穩固之手機支架+手機導航乃務實之選。但要注意：(1)手機屏幕小，駕駛時觀不清；(2)夏天暴曬將使手機過熱降頻甚至自動關機；(3)長時間導航耗電快，需配車載充電器。
  \item \textbf{建議}：短途或臨時用車，手機支架夠用；日常通勤長途，建議升級支持CarPlay/HiCar之車機或加裝後裝車機，體驗與安全皆更好。
\end{itemize}

\begin{tcbwarn}{安全提醒}
無論用哪種方案，駕駛時操作手機皆違反交規。手機支架只能用以“觀導航”，不可用以“操作手機”。設置目的地、切換歌曲等操作，請務必停車後再作，或使用語音控制。

\end{tcbwarn}

\booksection{六、實用術}

\begin{itemize}
  \item \textbf{無線CarPlay連不穩}：(1)關閉車機藍牙再重新配對；(2)刪除車機端已配對設備記錄後重連；(3)檢查手機為否同時連矣多個藍牙設備造成衝突；(4)升級車機固件與手機系統至最新版；(5)實於不行，改用有線CarPlay，更穩定。
  \item \textbf{車機流量用完矣}：(1)用手機開熱點，車機連手機熱點上網；(2)於車機設置中以"高耗流量"應用（如視頻、於線音樂）切至Wi-Fi；(3)定期清理車機離線地圖緩存。
  \item \textbf{駕駛模式自動觸發}：(1)iOS"專注模式"設定"連接車載藍牙時自動開啓"，自動屏蔽非緊急通知；(2)安卓"駕駛模式"或第三方App（如Tasker）可實現連藍牙自動靜音、自動回覆消息；(3)華爲手機可設置"智感支付"，上車自動切至車機播放音樂。
\end{itemize}

\begin{tcbinfo}{選車時之小建議}
若於乎車機體驗，選車時務必去4S店實測：(1)帶自己之手機試CarPlay/HiCar連接為否穩定；(2)試語音助手之反應速度與指令識別率；(3)觀應用市場有無汝常用之App；(4)關注後續OTA升級頻率——好車機乃用軟件迭代出來之。

\end{tcbinfo}

\begin{tcbpractice}{動手：優化汝之車機互聯體驗}
\begin{enumerate}
  \item 確認汝之手機與車機支持之互聯方案（CarPlay / HiCar / Android Auto / CarLife）。
  \item 清理一次車機之藍牙配對列表，唯保留常用手機，減少連接衝突。
  \item 設一個"駕駛專注模式"，連車機藍牙時自動開啓。
  \item 下載車機品牌App（如鴻蒙智行、理想汽車、蔚來App），觀為否能遠程控制空調、查看車況。
\end{enumerate}

\end{tcbpractice}

\begin{tcbsummary}{本章小結}
\begin{itemize}
  \item 車機互聯三大路線：投屏（CarPlay等）、深融合（鴻蒙座艙等）、CarPlay 2.0深集成
  \item CarPlay生態最廣；HiCar國內體驗好；Android Auto國內需折騰；CarLife適舊車後裝
  \item 鴻蒙座艙之超級桌面、應用接力、共用賬號代表矣深度融合方向
  \item 蔚來NOMI、理欲多屏、小米人車家、較亞迪旋轉屏各有特色
  \item 手機當車機僅適臨時場景，久宜升或後裝車機
  \item 實用術：CarPlay重連、手機熱點補流量、藍牙觸發駕駛模式
\end{itemize}

\end{tcbsummary}
